\clearpage\fancyhead[L]{\small PALLAS.}\hypertarget{native-359-lower}{}\label{native-359-lower}
\begin{center}\textbf{PALLAS.}\par\vspace{4pt}\hrule\vspace{4pt}\fontsize{10}{12}\selectfont Monatliche Correspondenz. Band XXVI. S. 199\;..\;203. 1812. August.\par\vspace{4pt}\hrule\vspace{4pt}\end{center}

\begin{flushright}Göttingen den 4. August 1812.\end{flushright}
\par\indent Ich habe Ihnen bereits vor einiger Zeit die erste Auffindung der \emph{Pallas} in diesem Jahre gemeldet.
Meiner eigenen Beobachtungen seitdem sind nur wenige. Um die Zeit der Opposition stand der
Planet die ganze Nacht hindurch für unser Locale zu hoch, um mit unserm besten Instrument micrometrische
Vergleichungen zu machen: am Mauerquadranten wurde er, da er nur zehnte Grösse hatte,
gar nicht sichtbar; ich habe mich also begnügt, ihn nur einmal mit einem schwächern Instrument um
diese Zeit am Kreismicrometer zu beobachten, wo dann die Declination wenig zuverlässig war.

\Needspace{11\baselineskip}\par\indent Meine sämmtlichen Beobachtungen sind folgende:

\begin{center}
\fontsize{10}{12}\selectfont\setlength{\tabcolsep}{3pt}\renewcommand{\arraystretch}{1.15}\begin{tabular}{c|c|c|c}
1812 & Mittlere Zeit & Gerade Aufsteig. & Nördl. Abw. \\
\hline
April 9 & $11^{\mathrm u}\;32^{\mathrm m}\;27^{\mathrm s}$ & $269^\circ\;5'\;58''4$ & $16^\circ\;38'\;0''\!:\,:$ \\
Mai 2 & 10 \quad 36 \quad 6 & 268 \quad 42 \quad 5.3 & 21 \quad 4 \quad 31.7 \\
4 & 10 \quad 44 \quad 44 & 268 \quad 31 \quad 38.0 & 21 \quad 25 \quad 4.3 \\
Juni 7 & 11 \quad 13 \quad 51 & 262 \quad 50 \quad 51.1 & 25 \quad 2 \quad 56.6 \\
\end{tabular}
\end{center}

\Needspace{9\baselineskip}\par\indent Dr. \textsc{Olbers} hatte in Bremen folgende zwei Beobachtungen gemacht:

\begin{center}
\fontsize{10}{12}\selectfont\setlength{\tabcolsep}{3pt}\renewcommand{\arraystretch}{1.15}\begin{tabular}{c|c|c|c}
1812 & Mittlere Zeit & Gerade Aufsteig. & Nördl. Abw. \\
\hline
April 3 & $11^{\mathrm u}\;26^{\mathrm m}\;44^{\mathrm s}$ & $268^\circ\;42'\;23''1$ & $15^\circ\;25'\;4''9$ \\
4 & 12 \quad 2 \quad 33 & 268 \quad 47 \quad 12.5 & 15 \quad 37 \quad 41.1 \\
\end{tabular}
\end{center}

\Needspace{12\baselineskip}\par\indent Von Paris aus hatte derselbe die Güte, mir folgende vier Meridianbeobachtungen des Dr. \textsc{Burkhardt}
mitzutheilen:

\begin{center}
\fontsize{10}{12}\selectfont\setlength{\tabcolsep}{3pt}\renewcommand{\arraystretch}{1.15}\begin{tabular}{c|c|c|c}
1812 & Mittlere Zeit & Gerade Aufsteig. & Nördl. Abw. \\
\hline
Juni 5 & $12^{\mathrm u}\;35^{\mathrm m}\;57^{\mathrm s}4$ & $263^\circ\;15'\;24''0$ & $24^\circ\;58'\;28''5$ \\
6 & 12 \quad 33 \quad 10.5 & 263 \quad 2 \quad 35. & 25 \quad 0 \quad 30. \\
8 & 12 \quad 21 \quad 37.0 & 262 \quad 37 \quad 6.5 & 25 \quad 3 \quad 57.7 \\
11 & 12 \quad 7 \quad 15.7 & 261 \quad 58 \quad 37.8 & 25 \quad 6 \quad 29.3 \\
\end{tabular}
\end{center}

\Needspace{20\baselineskip}\par\indent Die zweite Beobachtung war als weniger genau angegeben. Dieses sind alle Beobachtungen von
diesem Jahre, die mir bekannt geworden sind. Hr. \textsc{Nicolai} hat sie sämmtlich mit den Elementen verglichen
und folgende Übereinstimmung gefunden:
\SourceLeaf{360}\begin{center}
\emph{Unterschied.}
\end{center}

\begin{center}\fontsize{10}{12}\selectfont\setlength{\tabcolsep}{6pt}\renewcommand{\arraystretch}{1.15}\begin{tabular}{c|c|c|c}
1812 & AR. & Decl. & Beobachter\\\hline
April 3 & $+30''3$ & $-\!-10''7$ & Olbers\\
4 & $+33.3$ & $-16.8$ & Olbers\\
9 & $+30.6$ &  & Gauss\\
Mai 2 & $+57.6$ & $-5.4$ & Gauss\\
4 & $+49.4$ & $+6.6$ & Gauss\\
\end{tabular}\end{center}
\begin{center}\fontsize{10}{12}\selectfont\setlength{\tabcolsep}{6pt}\renewcommand{\arraystretch}{1.15}\begin{tabular}{c|c|c|c}
1812 & AR. & Decl. & Beobachter\\\hline
Juni 5 & $+51''4$ & $+6''1$ & Burckhardt\\
6 & $+59.4$ & $+11.7$ & Burckhardt\\
7 & $+56.7$ & $-34.0$ & Gauss\\
8 & $+59.6$ & $+2.4$ & Burckhardt\\
11 & $+62.0$ & $+8.2$ & Burckhardt\\
\end{tabular}\end{center}

\par\indent Dass die Elemente, bei welchen auf die Störungen durch den Jupiter auf das genaueste Rücksicht
genommen ist und die sich an alle sieben ersten Oppositionen so genau anschlossen, in der diesjährigen
Opposition 1 Minute abweichen, schreibe ich hauptsächlich dem Einfluss von Saturn und Mars
zu. Die vom erstern herrührenden Störungen hat Herr \textsc{Nicolai} unter meiner Aufsicht bereits einmal
vorläufig berechnet, und eine zweite schärfere Rechnung wird bald entweder von diesem ungemein geschickten
und fertigen Rechner oder von mir selbst vorgenommen werden, von deren Erfolg ich dann
zu seiner Zeit umständlicher sprechen werde.

\par\indent Die Opposition hat Herr \textsc{Nicolai} aus den Beobachtungen im Juni folgendermassen bestimmt.

\begin{center}\fontsize{10}{12}\selectfont\setlength{\tabcolsep}{0pt}\renewcommand{\arraystretch}{1.15}\begin{tabular}{l}
1812 Juni 10. $3^{\mathrm u}\;2^{\mathrm m}\;48^{\mathrm s}$\\Mittlere Zeit Meridian von Göttingen
\end{tabular}\end{center}
\begin{center}\fontsize{10}{12}\selectfont\setlength{\tabcolsep}{2pt}\renewcommand{\arraystretch}{1.15}\begin{tabular}{lcr}
Wahre Länge & $.\;.\;.\;.\;.\;.\;.\;.\;.\;.\;.\;.\;.\;.\;.$ & $259^\circ\;27'\;51''7$\\
Wahre geocentrische Breite & $.\;.\;.\;.\;.\;.\;.\;.\;.\;.$ & $48\;18\;26.1\;\mathrm{N}.$\\
\end{tabular}\end{center}

\Needspace{18\baselineskip}\par\indent Ich habe einstweilen versucht, wie genau sich sämmtliche acht bisher beobachtete Oppositionen
noch vereinigen lassen, wenn man blos auf die Störungen durch den Jupiter Rücksicht nimmt. Folgende
Verbesserungen an den Elementen müssen zu dem Behuf angebracht werden:

\begin{center}\fontsize{9}{11}\selectfont\setlength{\tabcolsep}{2pt}\renewcommand{\arraystretch}{1.15}\begin{tabular}{lcr}
Verbesserung der Epoche 1810 & $.\;.\;.\;.\;.\;.\;.\;.\;.\;.\;.$ & $-13''31$\\
der täglichen mittlern Bewegung & $.\;.\;.\;.\;.\;.\;.\;.\;.\;.$ & $-0.01200$\\
der Länge des Periheliums & $.\;.\;.\;.\;.\;.\;.\;.\;.\;.\;.\;.$ & $-17''57$\\
des Excentricitätswinkels & $.\;.\;.\;.\;.\;.\;.\;.\;.\;.\;.\;.\;.$ & $-4.09$\\
des Logarithmus der halben grossen Axe & $.\;.\;.\;.\;.\;.\;,$ & $+0.0000045$\\
der Länge des Knoten & $.\;.\;.\;.\;.\;.\;.\;.\;.\;.\;.\;.\;.\;.$ & $+0''80$\\
der Neigung der Bahn & $.\;.\;.\;.\;.\;.\;.\;.\;.\;.\;.\;.\;.\;.$ & $-1''21$\\
\end{tabular}\end{center}

\Needspace{16\baselineskip}\noindent Hiernach stimmen alle acht Oppositionen folgendermassen:

\begin{center}
\emph{Unterschied}
\end{center}

\begin{center}
\fontsize{10}{12}\selectfont\setlength{\tabcolsep}{6pt}\renewcommand{\arraystretch}{1.15}\begin{tabular}{c|c|c}
Opposition & Mittlere Länge & Helioc. Breite \\
\hline
1803 & $+5''8$ & $-2''8$ \\
1804 & $-12.9$ & $-9.7$ \\
1805 & $-15.2$ & $-0.7$ \\
1807 & $+8.3$ & $-5.8$ \\
1808 & $+16.0$ & $+11.7$ \\
1809 & $+13.0$ & $-1.5$ \\
1811 & $+11.7$ & $-9.5$ \\
1812 & $-25.9$ & $+6.1$ \\
\end{tabular}
\end{center}

\par\indent Ich bemerke hierbei nur noch, dass wo hier mitgetheilte Zahlen von den anderswo schon gedruckten
etwas differiren, dieses eine Folge einer später berichtigten Rechnung ist.

\Needspace{13\baselineskip}\par\indent Für das nächste Jahr hat Hr. \textsc{Nicolai} die Opposition wie folgt \emph{im Voraus} berechnet:
\SourceLeaf{361}\begin{center}\fontsize{10}{12}\selectfont\setlength{\tabcolsep}{0pt}\renewcommand{\arraystretch}{1.15}\begin{tabular}{l}
Zeit 1813 Aug. 18\quad $8^{\mathrm u}\;50^{\mathrm m}\;5^{\mathrm s}$\\Mittlere Zeit in Göttingen
\end{tabular}\end{center}
\begin{center}\fontsize{10}{12}\selectfont\setlength{\tabcolsep}{2pt}\renewcommand{\arraystretch}{1.15}\begin{tabular}{lcr}
Wahre Länge & $.\;.\;.\;.\;.\;.\;.\;.\;.$ & $325^\circ\;24'\;17''9$\\
Geocentr. Breite & $.\;.\;.\;,\;.\;.\;.\;.$ & $24\;37\;23.8\;\mathrm{N}.$\\
\end{tabular}\end{center}

\par\indent Die \emph{Juno}, welche dieses Jahr überaus lichtschwach war, habe ich auch mehreremale beobachtet;
von den daraus gefolgerten Resultaten werde ich in meinem nächsten Briefe eine umständliche Nachricht,
nebst einer von Herrn \textsc{Wachter} berechneten Ephemeride auf das nächste Jahr, mittheilen.
